% TOP_PROBLEM: {"id": 30007434, "title": "Rigidity of circle domains with removable boundary", "category_id": 8, "category": {"id": 8, "name": "analysis", "display_name": "Analysis", "description": "Limits, continuity, calculus, and function theory.", "slug": "analysis", "order_index": 8, "created_at": "2026-07-31T15:26:25.671Z"}, "status": "open"}
% FIRST_POSTED: 2026-10-06
% ATTEMPT_STATUS: solved
% OPENAI_CLAIM: {"schema_version":1,"source_repo":"https://github.com/openai/math","source_commit":"adc7f1241b42e322a6451854ab7e4b4c146bf78a","release_date":"2026-10-06","match":"full","resolution":"proved","independently_reviewed":false,"second_pass":"agreed","families":[{"family":"071","title":"Koebe’s circle-domain conjecture","subject":"Real and complex analysis","manuscripts":[{"title":"Removable Boundaries and Rigidity of Circle Domains","pdf_path":"preprints/Removable-Boundaries-and-Rigidity-of-Circle-Domains-September-23-2026/paper.pdf","readme_path":"preprints/Removable-Boundaries-and-Rigidity-of-Circle-Domains-September-23-2026/README.md","date":"2026-09-23","theorem_ref":"preprints/Removable-Boundaries-and-Rigidity-of-Circle-Domains-September-23-2026/build/sections/introduction.tex :: thm:main"}],"lean":{"doc_path":"lean/docs/071.md","comparators":["lean/ComparatorChallenges/KoebeCircleDomains.json","lean/ComparatorChallenges/KoebeCircleDomains.lean"],"declarations":[],"covers_main_theorem":true},"reasoning_trace":null}]}
\documentclass[11pt]{article}
\usepackage[margin=25mm]{geometry}
\usepackage{fontspec}
\setmainfont{Latin Modern Roman}
\usepackage{amsmath,amssymb,mathtools}
\usepackage{unicode-math}
\setmathfont{Latin Modern Math}
\usepackage{enumitem}
\usepackage{xurl,hyperref,bookmark}
\hypersetup{colorlinks=true,urlcolor=blue}
\setcounter{secnumdepth}{0}
\setlength{\parindent}{0pt}
\setlength{\parskip}{5pt}
\setlength{\emergencystretch}{3em}
\begin{document}
\section*{Rigidity of circle domains with removable boundary}

\subsection{OpenAI's claim}
\textit{Rigidity from removability}: The paper claims that a conformal map between circle domains must be a restriction of a Möbius map whenever the source boundary is conformally removable. No finiteness condition is placed on the collection of holes.

\subsection{Outline of the argument}
\textit{Rigidity from removability}: Scalar tests are carried into circle-domain models on successively larger compact regions. A Beltrami-based conductivity variation constructs low-energy coordinate barriers, which force any remaining omitted component to shrink to a point. The resulting scalar extension theorem and masks around disk components give a flux comparison; equality in the associated area estimates determines the conformal map.

\subsection{Scope relative to this problem}
\textit{Rigidity from removability}: The primary statement says: A circle domain with conformally removable boundary has the property that every conformal bijection onto another circle domain is a Mobius restriction. The target retains these quantified objects and hypotheses, so its positive assertion is covered.

With the hypotheses and conventions above, the claimed result addresses the whole mathematical assertion on this page.

\subsection{Formal verification}
Family 071 describes Lean coverage for the universal circle-domain representative and the Möbius rigidity implication under removable source boundary.

\textit{Rigidity from removability}: The manuscript links \texttt{KoebeCircleDomains.lean}.

The documented selection includes this manuscript’s main assertion; that assertion retains the restrictions stated above.

This site has not run the Lean development. These links describe OpenAI’s selected formal statements and should be read with their stated scope.

\subsection{Status caveat}
OpenAI reports that an internal model produced the manuscripts. Their arguments have not been independently verified as proofs by this site. Any portion outside the documented formal scope may contain mistakes, and the catalogue label reports the announced mathematical claim.

\subsection{Sources}
\begin{itemize}
\item OpenAI, \emph{Removable Boundaries and Rigidity of Circle Domains}, 2026-09-23; citation key \texttt{OAI:Removable-Boundaries-and-Rigidity-of-Circle-Domains-September-23-2026}. \href{https://github.com/openai/math/blob/adc7f1241b42e322a6451854ab7e4b4c146bf78a/preprints/Removable-Boundaries-and-Rigidity-of-Circle-Domains-September-23-2026/paper.pdf}{PDF (GitHub)}; \href{https://github.com/openai/math/raw/adc7f1241b42e322a6451854ab7e4b4c146bf78a/preprints/Removable-Boundaries-and-Rigidity-of-Circle-Domains-September-23-2026/paper.pdf}{PDF (direct)}; \href{https://github.com/openai/math/blob/adc7f1241b42e322a6451854ab7e4b4c146bf78a/preprints/Removable-Boundaries-and-Rigidity-of-Circle-Domains-September-23-2026/README.md}{Manuscript page and citation}; \href{https://github.com/openai/math/blob/main/preprints/Removable-Boundaries-and-Rigidity-of-Circle-Domains-September-23-2026/paper.pdf}{Latest version}.
\item Family 071: \href{https://github.com/openai/math/blob/adc7f1241b42e322a6451854ab7e4b4c146bf78a/lean/docs/071.md}{Lean scope}.
\item \href{https://github.com/openai/math/blob/adc7f1241b42e322a6451854ab7e4b4c146bf78a/lean/ComparatorChallenges/KoebeCircleDomains.json}{Comparator KoebeCircleDomains.json}.
\item \href{https://github.com/openai/math/blob/adc7f1241b42e322a6451854ab7e4b4c146bf78a/lean/ComparatorChallenges/KoebeCircleDomains.lean}{Comparator KoebeCircleDomains.lean}.
\item \href{https://github.com/openai/math/blob/adc7f1241b42e322a6451854ab7e4b4c146bf78a/CONTENTS.md}{OpenAI catalogue entry}.
\end{itemize}
\end{document}
